\newcommand{\DoNotLoadEpstopdf}{}
\documentclass[11pt,reqno]{amsart}
\usepackage[T1]{fontenc}
\usepackage{lmodern,amsmath,amssymb,amsthm,mathtools,mathrsfs,microtype}
\usepackage[margin=1.1in]{geometry}
\usepackage{xr-hyper}
\usepackage[hidelinks]{hyperref}
\externaldocument[supp-]{build/supplement}[supplement.pdf]
\hypersetup{pdftitle={Collision records, arithmetic, and raw local inversion in a triangular Lorentz gas},pdfauthor={Qian Qi}}
\setlength{\emergencystretch}{2em}
\allowdisplaybreaks[2]
\numberwithin{equation}{section}
\newtheorem{theorem}{Theorem}[section]
\newtheorem{lemma}[theorem]{Lemma}
\newtheorem{proposition}[theorem]{Proposition}
\newtheorem{corollary}[theorem]{Corollary}
\theoremstyle{definition}
\newtheorem{definition}[theorem]{Definition}
\newtheorem{inputstatement}[theorem]{Input statement}
\theoremstyle{remark}\newtheorem{remark}[theorem]{Remark}
\newcommand{\R}{\mathbb R}\newcommand{\Z}{\mathbb Z}
\newcommand{\one}{\mathbf1}\newcommand{\dd}{\,\mathrm d}
\newcommand{\esssup}{\operatorname*{ess\,sup}}
\newcommand{\diver}{\operatorname{div}}
\newcommand{\TV}{\operatorname{TV}}
\title[Collision records and raw local inversion]{Collision records, arithmetic, and raw local inversion in a triangular Lorentz gas}
\author{Qian Qi}
\date{October 12, 2026. A2-DYN, revision 85}
\subjclass[2020]{37D50, 37A50, 60F05, 28A75}
\keywords{Sinai billiards, exact collision records, local limits, maximal regular segments, residence-band height, microscopic concentration, coarea transport}
\begin{document}
\raggedbottom
\begin{abstract}
We retain the exact-label pointwise local limit problem for actual returns
in a triangular Lorentz gas. A Borel transversal parametrizes complete
maximal regular roof segments and identifies their angular disintegration
with the original collision measure, without a multiplicity factor.
Residence persistence then controls the essential height of every fixed
positive residence band, including critical and nongrazing boundary pairs.
Summing any finite collection of such bands yields a full surface-height
bound of order $r^{1/32}$ at every positive residence cutoff, with a leading
constant independent of that cutoff. The short critical part has the
sharper bound $r^{129/256}(1+|\log r|)$ in the same ordered norm.
The complete zero-residence tail remains inside the collision-normalized
height before the cutoff is removed. We identify it with the limiting
surface height and compare its actual boundary traces to the original
short-segment volume source by a paid fractional optical potential.
The original pointwise objective, exact labels and arithmetic classes
are unchanged. Its final budget retains two unproved microscopic
vanishing requirements: fractional optical concentration and
zero-residence concentration.
\end{abstract}
\maketitle

\section{The unchanged physical problem}
Put $R_-=9/20$ and $R_+=47/100$, and let $R\in I=[R_-,R_+]$. Circular scatterers have centers in
$\Lambda=\Z(1,0)+\Z(1/2,\sqrt3/2)$. Write $T_R$ for the regular
collision map and $(\alpha,p)\in(\R/2\pi\Z)\times(-1,1)$ for its
coordinates; $p$ is the outgoing tangential velocity and
$\dd\nu=\dd\alpha\,\dd p/(4\pi)$. Physical flights satisfy
$\tau_-\le\tau_R\le\tau_+$, where $\tau_-=3/50$ and $\tau_+$
is the fixed finite-horizon bound.

The section is specified here to fix every event in the argument.
Put $h_0=1/200$. Let $Y_R$ be the union of the four open coordinate
squares of half-side $h_0$ centered at
\[
 (0,0),\quad(\pi/6,-1/2),\quad
 (-\theta_R,\sin\theta_R),\quad
 (\pi/3-\theta_R,\sin\theta_R).
\]
Here $\theta_R=2\arctan q_R$, and $q_R$ is the unique root in
$(49/100,53/100)$ of
\[
 8Rq(1+q^2)-4q(1-q^2)+\sqrt3(1-6q^2+q^4)=0.
\]
Set
\[
 U=\{|\alpha-\pi/6|<3/50,\ |p|<3/20\},\qquad
 Y_R^*=Y_R\cup U,\qquad c_*=\frac{91}{10000\pi}.
\]
The rectangles are disjoint and $\nu(Y_R^*)=c_*$. We always use
$\nu_R^*=c_*^{-1}\nu|_{Y_R^*}$, not a renormalization of a smaller
source.

At $m$ flights, let $L=L_{m,R}$ be their total length,
$K^{\rm c}_{m,R}$ their lattice displacement, and
$A_{m,R}=\sum_{j=0}^{m-1}\one_{Y_R^*}\circ T_R^j$. The event is
\begin{equation}\label{eq:v76-event}
 \Pi_{n,k,m,R}=Y_R^*\cap\{K^{\rm c}_{m,R}=k,\ A_{m,R}=n\}
                         \cap T_R^{-m}Y_R^*.
\end{equation}
The initial visit is counted; the terminal visit is tested separately.
On this event the actual $n$th return has collision count $m$ and
roof $L$. Denote its density by $p_{n,R}(k,m,t)$.

Write $c_j=\sqrt{1-p_j^2}$ and let $\Delta_j$ be the least clearance
of flight $j$ from any third disk, excluding its two incident disks.
Keep the original smooth function $0\le\Theta\le1$, zero on
$(-\infty,1]$ and one on $[2,\infty)$. With $q_0=47/53$,
$d_j=\min(j,m-j)$ and $d'_j=\min(j,m-1-j)$, set
\begin{align}
 H^\varepsilon&=
 \prod_{j=0}^{m}\Theta\!\left(\frac{c_j}{\varepsilon q_0^{d_j/4}}\right)
 \prod_{j=0}^{m-1}\Theta\!\left(\frac{\Delta_j}{\varepsilon q_0^{d'_j/4}}\right),
 \qquad W^\varepsilon=1-H^\varepsilon,\label{eq:v76-guard}\\
 b^\varepsilon(t)\dd t
   &=L_\#(W^\varepsilon\one_{\Pi_{n,k,m,R}}\nu_R^*).
 \label{eq:v76-source}
\end{align}
In particular first incidence and first clearance outside any selected
critical neighborhoods are still present. No section-distance factor
has been inserted into this guard.

\subsection{The arithmetic reference and the target}
For $y=(k_1,k_2,t,n)$ put
\[
 \mu_R=(0,0,\bar\tau_R,c_*),\quad
 Z=(y-m\mu_R)/\sqrt m,\quad
 \mathcal C_{m,R,n,k}(M)=\{t:|Z|\le M\}.
\]
We retain the finite spectral-peak reference, including peaks close to
but not on the unit circle. To specify its meaning, the collision
record observable is $(\kappa_1,\kappa_2,\tau,\one_{Y_R^*})$.
Let $\lambda_j(R,\vartheta)$ and $\mathsf P_j(R,\vartheta)$ be its
continued simple twisted eigenvalue and projection in each of the
fixed peak charts, and let $\chi_j$ be their frequency partition.
At the real maximum $a_j$ of $\log|\lambda_j|$ put
\[
 z_j=\lambda_j(R,a_j),\quad
 \mu_j=i^{-1}\nabla\log\lambda_j(R,a_j),\quad
 H_j=-\nabla^2\log\lambda_j(R,a_j),
\]
with $\operatorname{Re}H_j\ge aI_4$ for a common $a>0$. If
$\eta_R=\one_{Y_R^*}$ and $\ell$ denotes mass evaluation, put
$B_j=\chi_j(R,a_j)\ell(M_{\eta_R}\mathsf P_j(R,a_j)(\eta_R\nu))$.
Define
\begin{align}
 G_H(z)&=(2\pi)^{-4}\int_{\R^4}e^{-iv\cdot z-v^{\mathsf T}Hv/2}\dd v,
 \notag\\
 \mathcal L_{m,R}(y)&=\operatorname{Re}\left[
 c_*^{-1}\sum_j e^{-ia_j\cdot y}z_j^m B_j
      G_{H_j}\left((y-m\mu_j)/\sqrt m\right)\right].
 \label{eq:v76-reference}
\end{align}
Only the charts active at $R$ are summed. This is the same transition
reference as in the retained spectral theorem. The existence and
uniform estimates for these scalar data belong to that analytic
input, not to the transport argument below. No frequency branch or
arithmetic zero class is deleted.

The governing assertion remains
\begin{equation}\label{eq:v76-target}
 \sup_{R,\,1\le n\le m,\,k\in\Z^2}
 \esssup_{t\in\mathcal C_{m,R,n,k}(M)}
 |m^2p_{n,R}(k,m,t)-\mathcal L_{m,R}(k_1,k_2,t,n)|\longrightarrow0.
\end{equation}
Let $\mathcal E_M$ denote this error's collision limsup. An empty
central target contributes zero. The order of the uniformity is
$\sup_R\esssup_t$, so exceptional density sets may depend on $R$.

\subsection{Complete surface estimates on all positive residence scales}
An actual nongrazing outgoing face carries its backward regular lifetime
$l=b-a$ in the same exact event. Let $\overline{\mathcal F}_r$ denote
its full first-exit assignment, averaged over roof increments in $(0,r)$.
A superscript $\ge\delta$ retains only faces with $l\ge\delta$.
Among its short faces, $l\le2r$, let $\mathcal K_r$ retain
$|g(b-)|^2<8\Lambda_\circ r$.

\begin{theorem}[All positive residence scales]
\label{thm:v85-main}
Under the explicit retained inputs below, for every sufficiently small
fixed $r>0$ and every fixed $\delta>0$,
\begin{align}
 \limsup_m\sup_{R,n,k}m^2
       \|\overline{\mathcal F}_r^{\ge\delta}\|_\infty
                 &\le C_\eta r^{1/32},\label{eq:v85-main-surface}\\
 \limsup_m\sup_{R,n,k}m^2\|\mathcal K_r^{\ge\delta}\|_\infty
                 &\le C_\eta r^{129/256}(1+|\log r|).
                                               \label{eq:v85-main-critical}
\end{align}
The constants are independent of the positive residence cutoff
$\delta$. No incoming-end type, physical word or exact label is removed
within these selected families. The same first bound holds for their
short noncritical paired incoming assignment.
\end{theorem}
Theorem~\ref{thm:v85-cutoff-surface} proves these assertions from
persistence on each complete residence band. The underlying height
bound is not inferred from local mass alone. Its measurable summation
is supplied by the explicit Borel bijection in
Theorem~\ref{thm:v85-segment-map}. Proposition~\ref{prop:v85-cutoff}
also reproduces the finite-count computation of the independent
first-inverse-gradient cutoff constant used for paired transfer.

The uniform leading constants do not authorize removal of $\delta$
before the collision limit. Theorem~\ref{thm:v85-boundary-charge}
identifies the precise remaining zero-residence charge and proves that
it is the small-width limit of the full surface height. Its comparison
with an actual collision-source restriction is paid by optical
variation, not by replacing a surface with a volume collar. Together
with the canonical fractional charge, it gives the two-term
microscopic budget in Theorem~\ref{thm:v85-endpoint}. Zero-charge
estimates remain required for the original pointwise theorem.

\section{Explicit analytic and geometric inputs}
We collect the indispensable inherited assertions in the exact form
used by the new proofs. Their source theorems and operator dependencies
are specified here and in Appendix~\ref{app:v83-kernels}. The new
potential and first-exit estimates do not assume a surface spectral
insertion theorem or an unweighted pointwise local bound.

\begin{inputstatement}[Physical flow and source comparison]
\label{input:v83-geometric}
On the regular part $\Omega$ of each exact event, endpoint angles
$q=(\alpha_0,\alpha_m)$ give $F=L$, $g=\nabla_qF=R(-p_0,p_m)$,
$K=D_q^2F$ and $\vartheta=\arg g$. Put
\[
 \lambda_\circ=9\sqrt{19}/200,\qquad \Lambda_\circ=47/6,\qquad
 \omega=|\partial_{p_0}\alpha_m|^{-1},\qquad
 \chi=|\det D_q(F,\vartheta)|.
\]
The bounds $\lambda_\circ I\le K\le\Lambda_\circ I$ and
$\lambda_\circ\le\chi\le\Lambda_\circ$ hold uniformly.
The explicit field
\[
 X_q=\frac{K^{-1}g}{g^{\mathsf T}K^{-1}g}
\]
is stopped at true physical and section-decision boundaries, not at
chart seams. Its flow satisfies $XF=1$, $X\vartheta=0$ and
\begin{equation}\label{eq:v83-radial-input}
 2\lambda_\circ s\le |g(\Phi_sx)|^2-|g(x)|^2
                                      \le2\Lambda_\circ s.
\end{equation}
Its endpoint-curve length is at most $C\sqrt s$. For
$I_m=1+\sum_{i=1}^{m-1}q_0^{d_i}/c_i$ one has
$|X\omega|\le C\omega I_m/|g|$ on regular survivors. At a genuine
grazing exit some internal incidence tends to zero and $\omega$ tends
to zero. Forward interior critical termination cannot occur.

Every starting point with forward lifetime at most $r$ belongs to a
layer of one of the following forms, at an actual collision or flight:
\begin{equation}\label{eq:v83-layer-input}
 c_j\le C_Eq_0^{d_j}\sqrt r,\quad
 \Delta_j\le C_Eq_0^{d'_j}\sqrt r,\quad
 \operatorname{dist}(T_R^jx,\partial Y_R^*)\le C_Eq_0^{d_j}\sqrt r.
\end{equation}
The physical coarea covariance, proved in the appendix in the convention
used here, has Jacobian $J_h=(\omega(\Phi_hx)/\omega(x))\chi(x)/\chi(\Phi_hx)$.
Writing $E_h^\varepsilon$ for all weighted starting exit mass and
\[
 O_h^\varepsilon(t)=\int_{S_h(t)} W^\varepsilon(x)
            [1-\omega(\Phi_hx)/\omega(x)]_+\,d\mu_t(x),
\]
the complete source satisfies, for each fixed $r$ and almost every $t$,
\begin{equation}\label{eq:v83-shift-input}
 b^\varepsilon(t)\le2A_rb^\varepsilon(t)+\bar E_r^\varepsilon(t)
       +\bar O_r^\varepsilon(t)
       +C_\Theta(1+\varepsilon^{-1})\sqrt r\,A_rp(t),
 \qquad A_rf(t)=\frac1{2r}\int_{t-r}^{t+r}f(s)\,ds.
\end{equation}
\end{inputstatement}
These are the endpoint-field, all-contact motion, uniform cubic-action
and separated-volume results in retained modules 174, 177, 180--182.
The guard term in \eqref{eq:v83-shift-input} is the complete product
guard, not a protective subword source. The derivative and continuation
assertions are confined to regular physical trajectories.

\begin{inputstatement}[Complete local source estimates]
\label{input:v83-local}
For every fixed auxiliary band $B_*$, constants $C_{B_*}<\infty$ and
$\rho_{B_*}<1$ give the following bounds uniformly in radius, exact
labels, interval $J$ and mark $0\le j\le m$. Let $S_{R,s}$ be the
union of strips of width $Cs$ about the finitely many coordinate lines
supporting $\partial Y_R^*$. For $s$ small,
\begin{align}
 m^2\nu\{K_m^{\rm c}=k,A_m=n,L_m\in J,
                       T_R^jx\in S_{R,s}\}
 &\le Cs^{1/16}(|J|+B_*^{-1})
   +C_{B_*}(1+B_*|J|)m^2\rho_{B_*}^{m/2},\label{eq:v83-strip-input}\\
 m^2\int_Jb^\delta(t)\,dt
 &\le C\delta^{1/16}(|J|+B_*^{-1})
   +C_{B_*}(1+B_*|J|)m^3\rho_{B_*}^{m/2}.\label{eq:v83-guard-input}
\end{align}
The local extension to a fixed remaining compact range of $\delta$
is included. At auxiliary parameter one $W^1=1$, so $b^1=p$.

Put $z=|g|$, $\alpha=129/128$ and, for $0<\eta<1/2$, define the
actual nonnegative source measures
\begin{align}
 v^{\eta,b}(t)\,dt&=L_\#(I_m z^{-\eta}\one_{\{0<z<b\}}
                                       \one_\Pi\nu_R^*),\notag\\
 v^\eta(t)\,dt&=L_\#(I_m z^{-\eta}\one_\Pi\nu_R^*).
                                                   \label{eq:v83-fractional-input}
\end{align}
They are finite and absolutely continuous at every fixed count, and
for $0<b<b_1$,
\begin{equation}\label{eq:v83-fractional-finite-input}
 m^2\int_Jv^{\eta,b}(t)\,dt
 \le C_\eta b^{\alpha-\eta}(|J|+B_*^{-1})
  +C_{B_*,\eta}(1+B_*|J|)m^2b^{1/2-\eta}\rho_{B_*}^{m/4}.
\end{equation}
Adding the bounded-gradient part gives, for every fixed $L>0$,
\begin{equation}\label{eq:v83-fixed-width-input}
 \limsup_m\sup_{R,n,k}\sup_{|J|=L}m^2\int_Jv^\eta\le C_\eta L,
 \qquad
 \limsup_m\sup_{R,n,k}\sup_{|J|=L}m^2\int_Jv^{\eta,b}
                                      \le C_\eta b^{\alpha-\eta}L.
\end{equation}
Every auxiliary band is fixed before the collision limit and enlarged
only afterwards; leading constants are independent of that band and
of the previously fixed interval length.
\end{inputstatement}
Equation~\eqref{eq:v83-strip-input} is Supplement
Theorem~\ref{supp-thm:v48-marked-strip-local}. The physical source
bound \eqref{eq:v83-guard-input} is Supplement
Theorem~\ref{supp-thm:v49-physical-remainder-local}, in particular its
finite-window estimate. The fractional bounds are the positive
layer-cake conclusions of retained module 187, stated in full above.
The joint spectral and common-space assumptions used there are recorded
in the appendix. No estimate for an untruncated first inverse gradient
or for a singular roof convolution is part of this input.

\begin{inputstatement}[The scalar positive remainder]
\label{input:v83-scalar}
For large fixed $B$, with $\varepsilon(B)=A_0B^{-1/12}$,
\begin{equation}\label{eq:v76-remainder}
 m^2p-\mathcal L=m^2b^{\varepsilon(B)}+e_B,\qquad
 \limsup_m\sup_{R,n,k}\esssup_{\mathcal C(M)}|e_B|
                                      \le C_MB^{-1/192}.
\end{equation}
The source $b^\varepsilon$ is nonnegative and has a compactly supported
$BV$ density at fixed count. No collision-uniform variation bound is
assumed. Supplement Theorem~\ref{supp-thm:v51-positive-error} gives
this remainder with the same complete finite-peak reference.
\end{inputstatement}

\input{core/190_untruncated_optical_potentials}
\input{core/191_complete_first_exit_surface_mass}
\input{core/192_ordered_fractional_tail_and_endpoint}

\input{core/196_measurable_segments_and_cutoff_audit}
\input{core/197_complete_residence_band_heights}
\input{core/198_zero_residence_charge_and_original_source}


\section{The original scalar objective and downstream normalization}
The governing assertion remains \eqref{eq:v76-target}. The full
surface is now estimated in essential height at every previously fixed
positive residence cutoff, with one leading constant for all such
cutoffs. The critical short surface has the stronger height estimate
\eqref{eq:v85-main-critical} under that same ordering. These are
complete positive-scale theorems, not a sum over selected words.

The complete zero-residence tail, and the complete fractional optical
tail, remain in the original normalized essential supremum before the
collision limit. Neither vanishing assertion follows by letting a
cutoff depend on collision count. Their exact role is
\eqref{eq:v85-endpoint}; the surface-source comparison
\eqref{eq:v85-trace-volume} does not suppress a boundary flux or change
the guard. The manuscript does not mark the organizing theorem proved
until these microscopic estimates are established.

The active dyadic formula is the offset expression
\eqref{eq:v83-dyadic-tail}; its negative-shift domain is explicit in
\eqref{eq:v85-charge-domain}. The historical module 189 remains
qualified at its inclusion in the supplement, with original bytes
preserved. The supplement also contains a compact input-to-theorem,
equation and proof-page map (Supplement
page~\pageref{supp-v85:dependency-guide}).

After scalar closure, the positive path-remainder equality in
Supplement Theorem~\ref{supp-thm:v73-scalar-path-equality} supplies
the same-roof numerator. Ratios are formed only on positive arithmetic
reference classes. The actual-return clock and forward likelihood
retain their separately stated hypotheses. No surface estimate is
used to reverse this scalar-first logical order.

\appendix
\input{appendices/v83_common_kernels_and_dependencies}
\begingroup
\small
\enlargethispage{2\baselineskip}
\input{journal-references}
\endgroup
\end{document}
